\section{Análise da Qualidade dos Dados}
Nesta seção, apresenta-se a auditoria de qualidade realizada sobre a base de dados integrada do SIA/SUS, utilizando as dimensões de completude e consistência.

\subsection{Matriz de Completude e Inventário}
A base de dados apresentou um índice de completude superior a 99\% para os campos de identificação municipal e período. No entanto, identificou-se que variáveis financeiras apresentam registros com valor zero, o que requer atenção na modelagem estatística.

\subsection{Priorização de Problemas (Pareto)}
Aplicou-se o Diagrama de Pareto para identificar os "poucos vitais" que afetam a integridade do banco. Conforme a Figura \ref{fig:pareto}, os valores zerados representam o maior volume de inconsistências.

\begin{figure}[h]
    \centering
    \includegraphics[width=0.7\textwidth]{figures/pareto_qualidade.png}
    \caption{Diagrama de Pareto para Problemas de Qualidade.}
    \label{fig:pareto}
\end{figure}

\subsection{Matriz GUT e Plano de Ação}
Para priorizar as correções, utilizou-se a Matriz GUT. O problema de "Duplicidade de Carga" foi classificado com a maior prioridade devido à sua tendência de crescimento caso o pipeline de ETL não possua travas de integridade.